%! The Secant, Cosecant, and Cotangent Functions — Unit 3, Lesson 3.11 — Warm-Up
\documentclass[10pt]{article}
\usepackage{precalculus-article}
\usepackage{precalculus-boxes}
\newcommand{\TallMath}[1]{$\displaystyle #1\rule[-1.4em]{0pt}{3.2em}$}

\begin{document}

\pageheader{Unit 3, Lesson 3.11}{Warm-Up}
\namedateperiod

\vspace{0.12in}

\begin{notesbox}{Spiral Review}

\textbf{1.}\quad Evaluate each expression using your knowledge of the unit circle.

\vspace{4pt}
\renewcommand{\arraystretch}{1.6}
\begin{tabular}{@{}p{0.23\linewidth}p{0.23\linewidth}p{0.23\linewidth}p{0.23\linewidth}@{}}
(a)\;$\cos\!\left(\dfrac{\pi}{3}\right) = $ \blank{1.1cm}
& (b)\;$\sin(\pi) = $ \blank{1.1cm}
& (c)\;$\cos\!\left(\dfrac{3\pi}{2}\right) = $ \blank{1.1cm}
& (d)\;$\sin\!\left(\dfrac{\pi}{6}\right) = $ \blank{1.1cm}
\end{tabular}

\vspace{0.14in}
\textbf{2.}\quad The tangent function is defined as a quotient of sine and cosine.
Fill in the definition and evaluate.

\vspace{4pt}
$\tan\theta = \dfrac{\underline{\hspace{1.8cm}}}{\underline{\hspace{1.8cm}}}$
\hspace{1cm}
(a)\quad $\tan\!\left(\dfrac{\pi}{4}\right) = $ \blank{1.5cm}
\hspace{0.8cm}
(b)\quad $\tan\!\left(\dfrac{\pi}{6}\right) = $ \blank{1.8cm}

\vspace{0.14in}
\textbf{3.}\quad The tangent function has vertical asymptotes at certain $\theta$ values.

\vspace{4pt}
(a)\quad What must be true about $\cos\theta$ for $\tan\theta$ to be undefined?

\vspace{4pt}
\writeline

\vspace{4pt}
(b)\quad List the two values of $\theta$ in $[0, 2\pi)$ where $\tan\theta$ is undefined:

\vspace{4pt}
$\theta = $ \blank{2cm} \quad and \quad $\theta = $ \blank{2cm}

\vspace{0.14in}
\textbf{4.}\quad For any nonzero number $x$, the \textbf{reciprocal} of $x$ is $\dfrac{1}{x}$.

\vspace{4pt}
(a)\quad Find the reciprocal of $\cos\!\left(\dfrac{\pi}{3}\right)$.
  Use your answer from Question 1(a).

\vspace{4pt}
$\dfrac{1}{\cos\!\left(\frac{\pi}{3}\right)} = \dfrac{1}{\underline{\hspace{1cm}}} = $ \blank{1.8cm}

\vspace{4pt}
(b)\quad Can the reciprocal of a number ever equal $\dfrac{1}{2}$? Explain why or why not.

\vspace{4pt}
\writeline

\end{notesbox}

\end{document}
